\section{Evolutionary Rationality: Formal Boundary Calculations}
\label{sec:annex-p19-rationality-1-3}

\begin{annexitem}{annex:p19.r13.success}{thm:p19-r13-success}{The task-count identity and its domain}
\begin{proof}
For \(a>0\), cancel the nonzero denominator in
\((c/a)(a/N)=c/N\). The identity counts correct rendered answers over
the complete equally weighted problem list. If \(a=0\), also \(c=0\);
the total proportion is zero but the conditional correctness fraction
\(c/a\) is undefined. Normalizing percentages to proportions is necessary
before multiplication. If problems have different costs or probabilities,
their unweighted count cannot recover the weighted loss. For instance one
error costing one and one error costing one hundred contribute the same
count but different total costs. The count identity neither specifies the
task distribution nor derives a norm or an evolutionary optimality claim.
\end{proof}\end{annexitem}

\begin{annexitem}{annex:p19.r13.stability}{thm:p19-r13-stability}{A local invasion calculation and its attraction basin}
\begin{proof}
Against a population playing stag with frequency \(x\), a stag player
has expected payoff \(5x+0(1-x)=5x\), while hare has constant payoff
three. Multiplying the payoff difference by \(x(1-x)\) gives the stated
two-strategy replicator equation. Its sign is negative on \((0,3/5)\)
and positive on \((3/5,1)\), so its interior trajectories move toward
the corresponding endpoint. The intermediate rest point is unstable.

For a pure stag resident perturbed by a fraction \(\epsilon\) hare,
resident payoff is \(5(1-\epsilon)\), mutant payoff three. The former
is larger for \(0<\epsilon<2/5\). More generally a mutant mixing stag
with positive hare proportion \(q\) earns
\((1-q)5(1-\epsilon q)+3q\), while a stag resident earns
\(5(1-\epsilon q)\). Their difference is
\(q[5(1-\epsilon q)-3]>0\) for sufficiently small positive
\(\epsilon\). This is precisely resistance to every distinct rare
two-strategy mutant mixture. Thus pure cooperation is locally ESS even
though a sufficiently large population change moves into the hare basin.
The algebra establishes the illustrative model's conditional claim, not
that an observed human group follows replicator dynamics or these payoffs.
\end{proof}\end{annexitem}

\begin{annexitem}{annex:p19.r13.horizon}{thm:p19-r13-horizon}{A normalized long-run gap is not zero loss on each occasion}
\begin{proof}
An agent loss one on the first occasion and zero subsequently has
cumulative loss \(L_T=1\). A comparator with zero loss on every
occasion has \(L_T^*=0\). Their per-occasion cumulative difference
is \(1/T\), which converges to zero although the first loss is
nonzero. Hence the asymptotic formula does not logically entail
correctness at every finite time. This record tests the strength of a
convergence conclusion; it is not a proof or refutation of any particular
meta-induction algorithm's guarantee under its own full hypotheses.
\end{proof}\end{annexitem}

\begin{annexitem}{annex:p19.r13.probabilistic-mt}{thm:p19-r13-probabilistic-mt}{The exceptional cell and the reversed conditional}
\begin{proof}
The four proposed cell masses are nonnegative and total one. They give
\(P(A)=1\) and \(P(B^c)=\epsilon>0\), so both conditional ratios
are defined. Their values are
\((1-\epsilon)/1=1-\epsilon\) and \(0/\epsilon=0\).
Choosing \(\epsilon\) arbitrarily small makes the first as close to
one as desired while leaving the second zero. The source material
conditional would instead exclude the exceptional \(A\cap B^c\)
cell altogether; assuming that different premise would support material
contraposition. The example changes premise semantics explicitly rather
than alleging a counterexample to valid material modus tollens.
\end{proof}\end{annexitem}
